% from ega0-22-fr.tex
\oldpage[0\textsubscript{IV-1}]{183}

\subsection{Relèvement de la lissité formelle}
\label{subsection:0.22.1-fr}

\begin{theorem}[22.1.1]
\label{0.22.1.1-fr}
Soient $s:\Lambda\to A$, $u:A\to B$ deux homomorphismes d'anneaux,
$\mathfrak J$ un idéal de $B$, $C$ l'anneau quotient $B/\mathfrak J$. On
suppose que $C$ est une $\Lambda$-algèbre formellement lisse pour la topologie
discrète.
\begin{enumerate}
  \item[(i)] Supposons vérifiées les conditions suivantes~:
  \begin{enumerate}
    \item[$\alpha$)] $\mathfrak J/\mathfrak J^2$ est un $C$-module projectif.
    \item[$\beta$)] L'homomorphisme canonique
    $\mathbf S_C^*(\mathfrak J/\mathfrak J^2)\to
    \operatorname{gr}_{\mathfrak J}^*(B)$ (19.5.2) est bijectif.
    \item[$\gamma$)] L'homomorphisme caractéristique
    $\chi:\Upsilon_{C/B/\Lambda}\to\mathfrak J/\mathfrak J^2$ (20.6.20)
    est injectif.
    \item[$\delta$)] Le $C$-module
    $\Omega^1_{B/A}\otimes_BC$ est projectif.
  \end{enumerate}
  Alors $B$, muni de la topologie $\mathfrak J$-préadique, est une
  $\Lambda$-algèbre formellement lisse.

  \item[(ii)] Inversement, supposons que $B$ soit une $\Lambda$-algèbre
  formellement lisse pour la topologie $\mathfrak J$-préadique, et que $A$
  soit une $\Lambda$-algèbre formellement lisse pour la topologie discrète.
  Alors les conditions $\alpha$), $\beta$), $\gamma$), $\delta$) de (i) sont
  vérifiées.
\end{enumerate}
\end{theorem}

\begin{proof}
En vertu de la suite exacte (20.6.22.1) (qui est applicable, puisque
$B/\mathfrak J^2$ est une extension $\Lambda$-triviale de $C$ par
$\mathfrak J/\mathfrak J^2$, $C$ étant supposée être une $\Lambda$-algèbre
formellement lisse (19.4.4)), la condition $\gamma$) équivaut à
$\Upsilon^C_{B/A/\Lambda}=0$, ou encore à la suivante~:
\begin{enumerate}
  \item[$\gamma'$)] L'homomorphisme
  \[
    u_{B/A/\Lambda}\otimes 1_C:
    \Omega^1_{A/\Lambda}\otimes_A C
    \longrightarrow \Omega^1_{B/\Lambda}\otimes_B C
  \]
  est injectif.
\end{enumerate}

(i) Les conditions $\alpha$) et $\beta$) entraînent que $B$ est une
$\Lambda$-algèbre formellement lisse pour la topologie
$\mathfrak J$-préadique (19.5.4). Il revient donc au même de dire que $B$ est
(pour la topologie $\mathfrak J$-préadique) une $A$-algèbre formellement
lisse, ou une $A$-algèbre formellement lisse \emph{relativement à $\Lambda$}.
Pour voir que $B$ est une $A$-algèbre formellement lisse, il suffit donc, en
vertu de (20.7.2), de prouver que l'homomorphisme continu de $B$-modules
topologiques
\[
  u_{B/A/\Lambda}:\Omega^1_{A/\Lambda}\otimes_A B
  \longrightarrow\Omega^1_{B/\Lambda}
\]
est \emph{formellement inversible à gauche}. Or, puisque $B$ est une
$\Lambda$-algèbre formellement lisse, le $B$-module topologique
$\Omega^1_{B/\Lambda}$ est formellement projectif (20.4.9) et sa topologie
est déduite de celle de $B$ (20.4.5). En vertu de (19.1.9), il suffit donc,
pour que $u_{B/A/\Lambda}$ soit formellement inversible à gauche, que
$u_{B/A/\Lambda}\otimes 1_C$ soit inversible à gauche. Mais en vertu de
$\gamma'$), cette dernière application est injective, donc on a la suite
exacte (20.6.14.7)
\[
  0\longrightarrow\Omega^1_{A/\Lambda}\otimes_A C
  \longrightarrow\Omega^1_{B/\Lambda}\otimes_B C
  \longrightarrow\Omega^1_{B/A}\otimes_B C\longrightarrow0.
\]
Enfin, en vertu de $\delta$), le $C$-module
$\Omega^1_{B/A}\otimes_BC$ est projectif, donc la suite exacte précédente
est scindée, ce qui achève de prouver (i).

